\begin{recipe}{Pasta amatriciana}
\kicker[Hoofdgerecht,Vlees,Italiaans,Pittig]{Hoofdgerecht · vlees · Italiaans · pittig}
\index[register]{Pasta amatriciana}
\index[register]{Guanciale}
\index[register]{Pecorino}
\index[register]{Tomatensaus}

\meta{45 minuten \dvd Voor 4 personen}

\lettrine{A}{matriciana} is een van de vier klassieke Romeinse pastagerechten,
samen met cacio e pepe, carbonara en alla gricia. De kinderen noemen het de rode carbonara, wat misschien niet zo gek is.

\blockrule{Ingrediënten}
\begin{ingredients}
  \ing{115 g}{guanciale of zuurkoolspek, in blokjes van 1\,cm}\index[register]{Guanciale}
  \ing{350 g}{spaghetti of bucatini}
  \ing{2 blikken (800\,g)}{gepelde tomaten, geplet}\index[register]{Tomaten}
  \ing{2 el}{olijfolie, plus extra}
  \ing{½ tl}{chilivlokken, plus extra naar smaak}
  \ingb{scheut Pinot grigio}
  \ingb{Pecorino Romano, geraspt}\index[register]{Pecorino}
  \ingb{zout en versgemalen peper}
\end{ingredients}

\blockrule{Bereiding}
\tip{Het gaat bij de guanciale om het uitbakken van het vet, niet om een krokant korstje. Neem er de tijd voor op laag vuur.}
\begin{steps}
  \item Breng een grote pan gezouten water aan de kook. 
  \item Verhit een grote hapjespan op laag vuur. Voeg het spek toe en bak,
        af en toe roerend, tot het vet grotendeels is uitgebakken en het spek
        lichtgoudbruin is, 10--15 minuten.
  \item Zet het vuur hoger. Voeg de chilivlokken toe en laat 30--60 seconden
        meebakken in het vet. Blus af met een scheut Pinot grigio en laat even
        koken.
  \item Voeg de tomaten toe en zet het vuur laag. Laat de saus 15--20 minuten
        zachtjes pruttelen tot hij indikt. Zet het vuur uit.
  \item Kook de pasta al dente: meestal een minuutje korter dan de verpakking
        aangeeft. Schep een flinke kop pastawater op voordat je afgiet.
  \item Voeg de pasta en een halve kop pastawater toe aan de saus. Kook op
        middelhoog vuur, al roerend, 3--4 minuten tot de saus mooi om de pasta
        heen zit. Voeg extra pastawater toe als je het iets dunner wilt.
        Haal van het vuur. Breng op smaak met zout, peper.
  \item Meng de pecorino er pas van het vuur af doorheen, of rasp hem alleen
        aan tafel: in een hete saus schift pecorino en wordt hij draderig.
  \item Verdeel over borden. Geef eventueel nog een scheutje olijfolie erbij.
        Wat plakjes komkommer of een groene salade erbij werkt goed als fris
        contrast.
\end{steps}
\heroimagefade[fadewidth=2.5cm]{images/pasta_amatriciana}
\end{recipe}
